\documentclass{article}

\usepackage{microtype}
\usepackage{graphicx}
\usepackage{subcaption}
\usepackage{booktabs}
\usepackage{tabularx}
\usepackage{hyperref}
\usepackage{amsmath}
\usepackage{amssymb}
\usepackage{multirow}
\usepackage{array}
\usepackage{xcolor}

\newcommand{\theHalgorithm}{\arabic{algorithm}}
\newcommand{\norm}[1]{\left\lVert #1 \right\rVert}
\newcommand{\R}{\mathbb{R}}
\newcommand{\od}{\mathbin{\odot}}
\newcommand{\DAWN}{DAWN-SRW}
\newcommand{\op}{\mathcal{O}}
\newcommand{\smallhead}[1]{\vspace{0.25em}\noindent\textbf{#1}}
\newcolumntype{Y}{>{\raggedright\arraybackslash}X}

\usepackage[preprint]{icml2026}
\makeatletter
\renewcommand{\Notice@String}{}
\makeatother
\hypersetup{pdfsubject={}}

\icmltitlerunning{\smash{DAWN-SRW: Learned State Transitions}}

\begin{document}

\twocolumn[
\icmltitle{DAWN-SRW: A Computational Model of Learned State Transitions}

\begin{icmlauthorlist}
\icmlauthor{Seungho Choi}{ind}
\end{icmlauthorlist}

\icmlaffiliation{ind}{Independent Researcher}
\icmlcorrespondingauthor{Seungho Choi}{madst0614@gmail.com}
\icmlkeywords{Language Modeling, Learned State Transitions, Read-Write Operators, Conditional Computation, Mechanistic Analysis}

\vskip 0.3in
]

\printAffiliationsAndNotice{}

\begin{abstract}
Language-model blocks evolve representation states through relational interaction and subsequent transformation. DAWN-SRW asks whether these otherwise implicit mechanisms can instead be constructed explicitly from reusable read-write (RW) primitives under state-conditioned control. It tests this architectural hypothesis by composing rank-1 RW primitives according to explicit state-conditioned applicability. Each current representation state determines which primitives participate and to what continuous degree. Their normalized composition constructs the query, key, and value representations used by attention, or directly produces the post-attention residual update. The goal is not to assume that learned operators are semantically human-readable, but to make transition-producing computation architecturally explicit enough to identify and causally intervene on its constituent operations. In the evaluated v4172 realization, learned operation-address mappings implement this control by comparing state-derived queries with bilinear addresses of primitive read/write directions. At 393.8M parameters and 39.999B tokens, DAWN-SRW retains performance against a matched dense Transformer: C4 cross-entropies are 2.9402 versus 2.9359, and six-task zero-shot means are 0.4629 versus 0.4597. A descriptive 1.272B evaluation reaches 2.8359 cross-entropy and a 0.4858 zero-shot mean. A discovery-frozen 4,096-site Q/K-centered operator circuit passes suppression, four matched-control families, and exact restoration on independent validation and held-out indirect-object-identification inputs. The model also exhibits exact input-conditioned functional support, although execution remains dense. These results support explicit state-conditioned RW composition as viable and causally analyzable, without claiming benchmark superiority or hardware efficiency.
\end{abstract}

\section{Introduction}
\label{sec:introduction}

Language-model blocks repeatedly transform representation states across depth.
Let \(X_\ell\) denote the sequence representation at layer \(\ell\); at a
coarse level, a block implements
\[
X_\ell\longrightarrow X_{\ell+1}.
\]
We use this state-evolution view only as an organizing observation, not as a
contribution.

In conventional Transformers, the mechanisms producing these state changes
remain largely implicit in dense parameterized transformations. Attention and
feed-forward layers expose module boundaries, but the architecture does not
directly represent which reusable elementary transformations apply to the
current state, under what conditions they apply, or which state-change
directions they produce.

Before asking whether learned neural computation is semantically interpretable,
we ask a more structural question: can the mechanism that transforms one
representation state into the next first be made explicit? Ideally, such a
representation would expose which reusable computational primitives apply to
the current state, to what degree they apply, how their contributions compose,
and how the resulting computation participates in the subsequent state update.
This does not require individual primitives to correspond to human-readable
concepts; it requires the transition-producing computation itself to be
architecturally identifiable.

This motivates the following architectural question:
\begin{quote}
\emph{Can the mechanisms producing language-model state evolution be
constructed explicitly from reusable RW primitives under state-conditioned
control?}
\end{quote}

We call this a reverse architectural hypothesis: instead of recovering
transition structure from an already-trained dense computation, we make that
structure explicit in the architecture and test whether it can support
end-to-end language modeling.

\DAWN{} investigates this question by taking a rank-1 read-write transformation
as its basic RW computational primitive. For a token state \(x\),
\begin{align}
\op_i(x)&=\langle x,r_i\rangle w_i,
\label{eq:intro_operator}\\
F(x)&=\frac{1}{Z(x)}\sum_i c_i(x)\op_i(x),
\qquad c_i(x)\in[0,1].
\label{eq:intro_field}
\end{align}
The read direction \(r_i\) determines what is extracted from the current
representation, the write direction \(w_i\) determines the direction of the
primitive's output contribution, and \(c_i(x)\) determines whether and to what
degree that primitive applies. Thus, \DAWN{} parameterizes transition-producing
computation as explicit conditional composition of reusable RW primitives.

The control abstraction is also explicit:
\begin{equation}
x\longmapsto C_a(x)
=\big(c_{a,1}(x),\ldots,c_{a,N_a}(x)\big).
\label{eq:intro_control}
\end{equation}
In v4172, this profile is realized through
\[
x\longmapsto\big(q_a(x),\tau_a(x)\big)
\longmapsto c_{a,i}(x),
\]
using primitive addresses \(k_{a,i}\). The operation-address space implements
this applicability control; it does not execute RW computation. Each RW
primitive defines a candidate computational contribution, while
state-conditioned applicability determines whether and to what degree that
contribution participates in the composed output.

This construction builds on the read-write view of dense feed-forward
computation, where first-layer directions detect input patterns and second-layer
directions contribute outputs \cite{geva2021transformer}. That analysis reveals
functional directions inside an existing dense layer. \DAWN{} instead turns RW
factorization into an architectural constraint: reusable RW units are
executable RW primitives whose state-dependent contributions are
directly controlled and observable.

The same formulation produces the query, key, and value states used by causal
self-attention and the direct update applied to the post-attention residual
state.
Q and K produce separate outputs from one shared QK pool; V and the
post-attention residual transition (RST) use separate pools. Operator pools and
applicability projections are reused across depth, while normalization and
attention output projections remain layer-specific.

The empirical study follows a fixed evidence order. We first compare a 393.8M
parameter \DAWN{} checkpoint with a parameter- and token-matched dense
Transformer on C4 and six zero-shot tasks. We then report descriptive 1.272B
results, measure exact input-conditioned operator support, and localize a
discovery-only indirect object identification (IOI) circuit. The circuit is
frozen before confirmatory data are opened, then tested with suppression,
matched controls, and exact restoration on independent validation and fully
held-out test inputs. Finally, we report three fail-closed boundaries: RAVEL
localization, ARC causal validation, and 1.3B IOI localization. Efficiency is
not a headline result: the current implementation scores and applies full
operator pools densely.

Our contributions are:
\begin{enumerate}
    \item \textbf{Explicit state-transition parameterization.} Starting from
    the view of language-model computation as representation-state evolution,
    we test whether its transition-producing transformations can be
    parameterized as compositions of explicit reusable RW primitives.
    \item \textbf{State-conditioned applicability.} We explicitly parameterize
    which RW primitives participate in the current computation and to what
    continuous degree. In the evaluated v4172 realization, this control
    additionally yields exact state-dependent support and is implemented
    through learned operation-address mappings.
    \item \textbf{End-to-end language modeling and causal analysis.} We show
    that this formulation can train language models at 400M and 1.3B scale while
    exposing exact operator support and causally intervenable operator
    contributions, including a discovery-frozen IOI circuit confirmed on
    independent validation and held-out inputs.
\end{enumerate}

\section{Related Work}
\label{sec:related}

\subsection{Conditional computation}

Sparsely-gated mixtures of experts activate a small subset of large expert
networks \cite{shazeer2017outrageously}, and Switch Transformers scale this
pattern with simple routing \cite{fedus2022switch}. PEER uses product-key
retrieval over more than a million small experts \cite{he2024mixture}. These
methods formulate conditional computation as expert retrieval; \DAWN{} instead
organizes the base language-model computation explicitly around RW primitives
and their state-dependent applicability. Its
current implementation still uses a dense address scan and dense, chunked RW
application; functional sparsity is not an indexed sparse kernel.

\subsection{Read-write primitives and memory views}

Transformer blocks combine attention and position-wise feed-forward networks
over a residual stream \cite{vaswani2017attention}. Geva et al. interpret the
feed-forward layers as key-value memories, with read-side pattern detection and
write-side output directions \cite{geva2021transformer}. Product Key Memory
factorizes retrieval keys \cite{lample2019large}, and later work scales
trainable key-value memory layers \cite{berges2024memory}. MoRAM further treats
rank-1 read/write factors as atomic associative-memory experts and uses the
read-side response itself to determine relevance \cite{lu2026little}.
\DAWN{} differs in using RW primitives as the language model's base
computational substrate and evaluating applicability through a separate
learned address representation.

\subsection{Mechanistic analysis and causal intervention}

Mechanistic interpretability studies model behavior through identifiable
components and causal interventions. Work on IOI isolates an attention-head
circuit and evaluates faithfulness, completeness, and minimality
\cite{wang2023interpretability}. Activation patching and scalable
approximations localize behavior to components \cite{kramar2024atp}. RAVEL
emphasizes that neuron-level representations may be entangled and evaluates
methods against causal criteria \cite{huang2024ravel}. Our analysis treats a
site as a layer, computation type, and RW-operator identity; freezes
localization before confirmation; and uses direct numerator suppression,
matched random circuits, and exact same-example restoration. Explicit RW units
make the intervention target precise, but do not guarantee a human-readable
interpretation.

\section{DAWN-SRW}
\label{sec:method}

\subsection{State Evolution as the Starting Abstraction}
\label{sec:state-evolution}

Let
\[
X_\ell=(x_{\ell,1},\ldots,x_{\ell,T})\in\R^{T\times D}
\]
denote the sequence representation at depth \(\ell\). At a coarse level, a
Transformer-like block evolves this state through a relation-dependent
attention update followed by a transformation of the contextualized
representation:
\begin{align}
Y_\ell&=X_\ell+A_\ell(X_\ell),\\
X_{\ell+1}&=Y_\ell+R_\ell(Y_\ell).
\label{eq:state_evolution}
\end{align}
This decomposition only fixes the computational abstraction and notation.
\DAWN{} concerns how the transformations producing these updates are
parameterized.

\subsection{Explicit State-Transition Hypothesis}
\label{sec:transition-hypothesis}

\DAWN{} tests the reverse hypothesis that transition-producing computation can
be constructed from explicit reusable RW primitives under explicit
state-dependent control. For a token state \(x\in\mathcal{S}=\R^D\), let
\(a\in\{Q,K,V,\mathrm{RST}\}\) index the computation type and define
\begin{equation}
F_a:\mathcal{S}\rightarrow\R^D.
\label{eq:composed_output}
\end{equation}
The Q, K, and V outputs form the state-conditioned interfaces through which
relational interaction is computed, whereas \(F_{\mathrm{RST}}\) directly
specifies a residual-state update. Each composed output is explicitly
decomposed into contributions from reusable RW primitives under state-dependent
control. The construction therefore separates the primitives, their
applicability profile, and their normalized composition. The operation-address
space is the evaluated implementation of this control mechanism rather than the
paper's core hypothesis.

\begin{figure*}[t]
\centering
\setlength{\fboxsep}{4pt}
\small
\textbf{Language modeling as representation-state evolution}\\[0.35em]
\begin{tabular}{c@{\hspace{0.12em}\(\longrightarrow\)\hspace{0.12em}}c@{\hspace{0.12em}\(\longrightarrow\)\hspace{0.12em}}c@{\hspace{0.12em}\(\longrightarrow\)\hspace{0.12em}}c@{\hspace{0.12em}\(\longrightarrow\)\hspace{0.12em}}c}
\fbox{\parbox[c][0.92cm][c]{0.10\textwidth}{\centering
\textbf{Sequence state}\\\(X_\ell\)}} &
\fbox{\parbox[c][0.92cm][c]{0.17\textwidth}{\centering
\textbf{Relational interaction}\\attention via \(F_Q,F_K,F_V\)}} &
\fbox{\parbox[c][0.92cm][c]{0.10\textwidth}{\centering
\textbf{Contextual state}\\\(Y_\ell\)}} &
\fbox{\parbox[c][0.92cm][c]{0.16\textwidth}{\centering
\textbf{Residual-state transition}\\via \(F_{\mathrm{RST}}\)}} &
\fbox{\parbox[c][0.92cm][c]{0.10\textwidth}{\centering
\textbf{Next state}\\\(X_{\ell+1}\)}}
\end{tabular}\\[-0.05em]
\begin{tabular}{cc}
\parbox{0.34\textwidth}{\centering
\(F_Q,F_K,F_V\) parameterize relational interaction\\\(\Downarrow\)} &
\parbox{0.34\textwidth}{\centering
\(F_{\mathrm{RST}}\) directly updates the residual state\\\(\Downarrow\)}
\end{tabular}\\[0.15em]
\textbf{\DAWN{}: explicit construction of state-conditioned computation}\\[0.35em]
\begin{tabular}{c@{\hspace{0.12em}\(\longrightarrow\)\hspace{0.12em}}c@{\hspace{0.12em}\(\longrightarrow\)\hspace{0.12em}}c@{\hspace{0.12em}\(\longrightarrow\)\hspace{0.12em}}c}
\fbox{\parbox[c][0.92cm][c]{0.09\textwidth}{\centering
\textbf{Token state}\\\(x\)}} &
\fbox{\parbox[c][0.92cm][c]{0.20\textwidth}{\centering
\textbf{State-conditioned}\\\textbf{applicability}\\\(C_a(x)\)}} &
\fbox{\parbox[c][0.92cm][c]{0.18\textwidth}{\centering
\textbf{RW primitive}\\\(\op_{a,i}(x)=\langle x,\bar r_{a,i}\rangle\bar w_{a,i}\)}} &
\fbox{\parbox[c][0.92cm][c]{0.13\textwidth}{\centering
\textbf{Composed output}\\\(F_a(x)\)}}
\end{tabular}\\[0.18em]
\parbox{0.82\textwidth}{\centering\footnotesize
\(C_a(x)=(c_{a,1}(x),\ldots,c_{a,N_a}(x))\); v4172:
\(q_a(x),\tau_a(x),k_{a,i}\) determine \(c_{a,i}(x)\)}
\caption{The two-level construction. Top: a language-model block evolves a
sequence representation through relational interaction and a subsequent
residual-state transition. Bottom: \DAWN{} constructs each composed output from
explicit RW primitives under state-conditioned applicability. In v4172, the
operation-address mapping implements the applicability step; RW execution
remains in the full \(D\)-dimensional representation space.}
\label{fig:architecture}
\end{figure*}

\subsection{Read-Write Computational Primitives}
\label{sec:rw}

Each primitive owns learned read and write parameters. With
\(\operatorname{unit}(z)=z/\max(\norm{z}_2,\epsilon)\), define normalized
directions \(\bar r_{a,i}=\operatorname{unit}(r_{a,i})\) and
\(\bar w_{a,i}=\operatorname{unit}(w_{a,i})\). The primitive executes
\begin{equation}
\op_{a,i}(x)=\langle x,\bar r_{a,i}\rangle\bar w_{a,i}
=\bar w_{a,i}\bar r_{a,i}^{\top}x.
\label{eq:rw_operator}
\end{equation}
This is an explicit rank-1 RW primitive. It reads a scalar from the current
state and produces a directed output contribution. The read direction
specifies the extracted scalar, while the write direction specifies the
direction of that contribution. For Q, K, and V, these contributions compose
the interfaces used by attention; for RST, they compose a direct residual-state
update. Unlike a static retrieved value, an RW primitive remains an
input-dependent transformation, with its output varying with the current state
through the read response. The state-to-state transition
is completed only when the composed output participates in the
attention-mediated or residual update.

\subsection{State-Conditioned Applicability}
\label{sec:applicability}

At the abstract level, control is represented by an applicability profile over
reusable RW primitives. For each computation type
\(a\in\{Q,K,V,\mathrm{RST}\}\), the current state induces a profile over all
\(N_a\) primitives,
\begin{equation}
C_a(x)=\big(c_{a,1}(x),\ldots,c_{a,N_a}(x)\big),
\label{eq:applicability_profile}
\end{equation}
whose coordinates specify which primitives participate and to what degree,
rather than assigning the state to a single discrete branch. The evaluated
v4172 realization has exact zero support outside learned boundaries and
continuously positive applicability inside them. This is not an algebraic
decomposition of the representation vector; its coordinates control the
contributions of the corresponding RW primitives composed below.

In v4172, normalized coordinates \(q_a(x)\) and \(k_{a,i}\) compare the current
state with primitive \(i\)'s learned address;
Section~\ref{sec:operation-space} gives their construction. A separate direct
projection produces the state-dependent DirectTau boundary:
\begin{align}
\tau_a(x)&=-1+2\sigma(xu_a+\beta_a),\\
\rho_{a,i}(x)&=q_a(x)^\top k_{a,i},\\
m_{a,i}(x)&=\rho_{a,i}(x)-\tau_a(x).
\label{eq:directtau}
\end{align}
The evaluated \texttt{linear\_angular} rule defines primitive applicability as
\begin{equation}
c_{a,i}(x)=
\operatorname{clip}\!\left(
\frac{m_{a,i}(x)}
{\max(1-\tau_a(x),10^{-4})},0,1\right).
\label{eq:angular_gate}
\end{equation}
We call \(c_{a,i}(x)\in[0,1]\) the \emph{continuous conditional
applicability} of primitive \(i\) at state \(x\). Its exact state-dependent
support is
\begin{equation}
\mathcal{I}_a(x)=\{i:\rho_{a,i}(x)>\tau_a(x)\}
=\{i:c_{a,i}(x)>0\}.
\label{eq:applicable_support}
\end{equation}
Outside this set, a primitive contributes exactly zero; inside it,
\(0<c_{a,i}(x)\leq1\). A primitive's contribution therefore enters
continuously from zero at the learned boundary. This combines state-dependent
sparse support with a continuous-valued degree of applicability rather than
merely assigning a dense soft mixture over every primitive.

\subsection{State-Conditioned Operator Composition}
\label{sec:operator-composition}

Applicable primitives compose into the output \(F_a(x)\):
\begin{align}
Z_a(x)&=\left[\max\!\left(\sum_i c_{a,i}(x),1\right)\right]^{p_a},\\
s_a&=\sqrt{D/L},\nonumber\\
F_a(x)&=\frac{s_a}{Z_a(x)}
\sum_i c_{a,i}(x)\op_{a,i}(x).
\label{eq:implemented_composition}
\end{align}
Equivalently,
\begin{equation}
\begin{aligned}
F_a(x)&=A_a(x)x,\\
A_a(x)&=\frac{s_a}{Z_a(x)}
\sum_i c_{a,i}(x)\bar w_{a,i}\bar r_{a,i}^{\top}.
\end{aligned}
\label{eq:state_dependent_operator}
\end{equation}
This form makes the architectural hypothesis explicit. The model does not apply
one fixed dense transformation to every state. Instead, the current
representation determines continuous applicability coefficients, which
construct a state-dependent transformation \(A_a(x)\) from reusable rank-1 RW
primitives. The resulting computation is both input-conditioned and explicitly
decomposed into identifiable operator contributions.

``Continuous'' refers to the value of operator applicability, not to
continuous-time dynamics or a continuous-depth model; layer evolution remains
discrete. The scalar mapping from normalized margin to applicability is
continuous and piecewise linear, with nondifferentiable support boundaries.

The denominator powers are \(p_Q=p_K=0.5\), \(p_V=1.0\), and
\(p_{\mathrm{RST}}=1.2\). The fixed composition scale is
\(s_a=\sqrt{D/L}\), with no
learned pool-strength parameter.

\subsection{Operation-Address Space}
\label{sec:operation-space}

The operation-address space \(\mathcal{A}=\R^R\) is where primitive
applicability is evaluated, not where RW computation itself is executed. This
addressing construction is one evaluated realization of the state-conditioned
applicability abstraction above. The present experiments do not establish that
this particular realization is necessary or preferable to alternative
state-conditioned control mechanisms. In v4172, \(R=256\), and a normalized
query is a direct projection of the current input state:
\begin{equation}
q_a(x)=\operatorname{unit}(xP_a+b_a)\in\mathcal{A}.
\label{eq:query}
\end{equation}
A shared probe pair \(P_r,P_w\in\R^{D\times R}\) maps all operator pools'
functional read and write directions into address features. The
operator address is
\begin{align}
u_{a,i}&=\operatorname{unit}(\bar r_{a,i}^{\top}P_r),&
v_{a,i}&=\operatorname{unit}(\bar w_{a,i}^{\top}P_w),\nonumber\\
k_{a,i}&=\operatorname{unit}(u_{a,i}\od v_{a,i})\in\mathcal{A}.
\label{eq:rw_address}
\end{align}
This generalized coordinate-wise bilinear mapping makes address identity a
learned function of what a primitive reads and where it writes, rather than an
independent expert ID. It is the addressing realization used to compute the
applicability coefficients in Equation~\ref{eq:angular_gate}, not the paper's
definition of a state transition. Shared probes align QK, V, and RST addresses
to one coordinate convention, while their operator rows remain in QK-, V-, and
RST-specific pools.

Because both \(q_a(x)\) and \(k_{a,i}\) are unit-normalized, they lie on the
address sphere \(\mathbb{S}^{R-1}\). Let
\(\theta_{a,i}(x)=\arccos\!\big(q_a(x)^\top k_{a,i}\big)\) and
\(\theta_a(x)=\arccos\tau_a(x)\). The active support is therefore
\[
\mathcal{I}_a(x)
=\{i:\theta_{a,i}(x)<\theta_a(x)\},
\]
a spherical cap centered at the query \(q_a(x)\) with state-dependent angular
radius \(\theta_a(x)\). Geometrically, \(q_a(x)\) specifies the center of the
applicability query, \(\tau_a(x)\) specifies its angular support boundary, and
\(c_{a,i}(x)\) specifies the continuous degree to which primitive \(i\)
applies. When the
numerical guard in Equation~\ref{eq:angular_gate} is inactive, an active
primitive has
\[
c_{a,i}(x)=
\frac{\cos\theta_{a,i}(x)-\tau_a(x)}{1-\tau_a(x)},
\]
which rises continuously from zero at the cap boundary to one at its center.
Equivalently, \(c_{a,i}(x)\) is a normalized continuous depth within the cap.
Equation~\ref{eq:angular_gate} is the exact implementation, retaining the
\(10^{-4}\) denominator guard and clipping.

\subsection{Attention-Mediated and Residual Transitions}
\label{sec:block}

For a sequence state \(X_\ell\), layer-specific normalization gives
\(\widetilde X_\ell=\operatorname{LN}^{(1)}_\ell(X_\ell)\). Three
state-conditioned operator compositions produce the attention interfaces:
\begin{align}
Q_\ell&=F_Q(\widetilde X_\ell),\nonumber\\
K_\ell&=F_K(\widetilde X_\ell),\nonumber\\
V_\ell&=F_V(\widetilde X_\ell),\nonumber\\
A_\ell&=\operatorname{softmax}\!\left(
\frac{Q_\ell K_\ell^\top}{\sqrt{d_h}}+M_{\mathrm{causal}}\right)
\!,\nonumber\\
\Delta X_{\mathrm{attn},\ell}&=A_\ell V_\ell W^O_\ell,\\
Y_\ell&=X_\ell+\Delta X_{\mathrm{attn},\ell}.
\label{eq:attention}
\end{align}
Q and K produce separate outputs from the same 4,600-row QK pool; V produces an
output from a separate 14,500-row pool. After attention, a second layer-specific
normalization \(\widetilde Y_\ell=\operatorname{LN}^{(2)}_\ell(Y_\ell)\) is
evaluated anew using the 29,300-row RST operator pool:
\begin{equation}
X_{\ell+1}=Y_\ell+F_{\mathrm{RST}}(\widetilde Y_\ell).
\label{eq:rst}
\end{equation}
The block thus applies an attention-mediated state transition followed by a
direct residual-state transition. RST applicability depends on the
contextualized post-attention state, not the pre-attention query.

\subsection{Cross-Depth Operator Reuse}
\label{sec:reuse}

The QK, V, and RST read/write pools; the two RW-to-address probes; the Q/K/V
query projection; the RST query projection; and their DirectTau projections
are shared across all \(L\) layers. Q and K share read/write rows but retain
separate query slices, tau values, applicability coefficients, and outputs. The
layer-specific parameters are two LayerNorms and the attention output
projection \(W^O_\ell\). Positional and token embeddings and the final
normalization are model-level parameters. Consequently, \DAWN{} reuses
RW primitives across depth without claiming complete block sharing.

\subsection{Implemented Addressing and Compute Boundary}
\label{sec:implemented-boundary}

The implementation shards operator rows across the model mesh and processes
them in dense chunks. The 400M configuration uses one QK chunk, one V chunk,
and four RST chunks. Every query scores the full local address pool; every
operator row participates in dense read and write contractions before
cross-shard reduction. Exact zeros in Equation~\ref{eq:angular_gate} therefore
describe \emph{functional} support, not an indexed sparse kernel.

The evaluated architecture also keeps RW execution in the full residual
dimension \(D\). The low-dimensional space is only an applicability-address
space. An executor could omit exactly zero RW terms without changing the
mathematical output, but that conditional execution opportunity is not
implemented or timed here.

\section{Experimental Setup}
\label{sec:setup}

\subsection{Models and training}

Table~\ref{tab:model-spec} summarizes the evaluated checkpoints. The primary
\DAWN{} model has 393,800,708 measured parameters, operation-address dimension
\(R=256\), and 48,400 operator rows across QK, V, and RST. It was trained for
39,999,504,384 packed C4 tokens with sequence length 512 and global training
batch 1,024. Both dropout rates are zero.

\begin{table*}[t]
\centering
\caption{Model and training specification. The 1.3B row is descriptive and
does not form a controlled scaling curve with the 400M rows.}
\label{tab:model-spec}
\small
\resizebox{0.99\textwidth}{!}{
\begin{tabular}{lrrrrrrrr}
\toprule
Model & Parameters & Tokens & \(D\) & Layers & Heads & QK / V / RST rows & Checkpoint\\
\midrule
\DAWN{} 400M & 393,800,708 & 39,999,504,384 & 2,304 & 18 & 36 &
4,600 / 14,500 / 29,300 & 76,293\\
Dense Transformer 400M & 393,896,960 & 39,999,504,384 & 1,280 & 18 & 20 &
-- & 76,293\\
\DAWN{} 1.3B & 1,272,038,404 & 19,999,981,568 & 4,096 & 24 & 32 &
8,534 / 26,904 / 54,352 & 87,193\\
\bottomrule
\end{tabular}
}
\end{table*}

\subsection{Matched 400M dense control}

The dense control has 393,896,960 measured parameters, width 1,280, feed-forward
width 5,120, 18 layers, and 20 heads. It uses the same tokenizer, vocabulary,
sequence length, training-token count, seed, batch size, mesh, learning rate,
warmup ratio, weight decay, and zero dropout as the primary model. The
parameter difference is 96,252, or 0.0244\% of the dense model. This comparison
tests performance retention under matched scale and training exposure; it is
not a formal equivalence test.

\subsection{Evaluation protocols}

The C4 protocol evaluates 610 complete global batches from the fixed
10M-token capped packed-validation view: 19,520 sequences and 9,974,720 valid
next-token targets. Cross-entropy, perplexity, and next-token accuracy use
float32 scoring. The zero-shot protocol uses \texttt{lm-eval} 0.4.2 with no
few-shot examples and no task limit on LAMBADA, HellaSwag, PIQA, ARC-Easy,
ARC-Challenge, and WinoGrande. The primary metric is raw accuracy for LAMBADA
and WinoGrande and normalized accuracy for the other four tasks. The mean is
an unweighted mean of these six metrics.

\subsection{Mechanistic localization and frozen intervention protocol}
\label{sec:protocol}

We define a site by layer, computation type, token position, and global operator
identity.
IOI discovery uses 119 independent paired-correct units. Sites are ranked by
absolute discovery mean contribution importance. We choose the smallest
audited prefix that reaches at least 0.70 cumulative absolute importance and
0.95 split top-\(k\) overlap. The resulting 4,096-site Q/K/V/RST/layer
composition and intervention specification are frozen before validation.

The primary causal score is the summed-log-probability margin between the
correct indirect object and the source name. Circuit-wide suppression zeros
only the selected operators' execution numerator at target positions while
retaining the full production denominator. Four control families each contain
100 frozen-size circuits that exclude circuit sites: equal-count random,
same-layer random, activation-matched, and route-frequency-matched across
Q/K/V/RST. Exact
restoration reinserts the same-example selected numerator after suppression,
again under the production denominator. Validation uses 123 independent
paired-correct units; the held-out test accessor is opened only after the
validation result and unchanged specification are finalized. We report 95\%
bootstrap confidence intervals and two-sided paired permutation tests where
specified.

\section{Results}
\label{sec:results}

\subsection{Language-model performance retention at 400M}

Table~\ref{tab:c4} reports the locked C4 evaluation. \DAWN{} differs from the
dense control by \(+0.004318918\) cross-entropy, \(+0.081537182\) perplexity,
and \(-0.002188432\) accuracy (\(-0.2188\) percentage point). These small
measured gaps support performance retention, not superiority or formal
equivalence.

\begin{table}[t]
\centering
\caption{Packed C4 validation on the identical 9,974,720 target tokens.}
\label{tab:c4}
\small
\begin{tabular}{lrrr}
\toprule
Model & CE & PPL & Accuracy\\
\midrule
\DAWN{} 400M & 2.940213 & 18.919873 & 0.424912\\
Dense 400M & 2.935894 & 18.838336 & 0.427100\\
\DAWN{} 1.3B & 2.835879 & 17.045382 & 0.437099\\
\bottomrule
\end{tabular}
\end{table}

Table~\ref{tab:zero-shot} gives task-level zero-shot results with standard
errors. The matched 400M means are 0.462872 and 0.459652, a difference of
\(+0.003220495\) (\(+0.3220\) percentage point). Task differences vary in
sign, so the mean does not establish benchmark superiority. LAMBADA
perplexities are \(24.0774\pm0.8005\) for \DAWN{} 400M,
\(24.0042\pm0.8073\) for dense 400M, and \(17.6358\pm0.5390\) for
\DAWN{} 1.3B.

\begin{table*}[t]
\centering
\caption{Zero-shot accuracy. Parentheses contain standard errors. HS, AE, and
AC denote HellaSwag, ARC-Easy, and ARC-Challenge.}
\label{tab:zero-shot}
\scriptsize
\resizebox{\textwidth}{!}{
\begin{tabular}{lrrrrrrr}
\toprule
Model & LAMBADA & HS & PIQA & AE & AC & WinoGrande & Mean\\
\midrule
\DAWN{} 400M &
0.352028 (0.006654) & 0.458375 (0.004972) & 0.712187 (0.010563) &
0.484007 (0.010255) & 0.256826 (0.012767) & 0.513812 (0.014047) & 0.462872\\
Dense 400M &
0.353192 (0.006659) & 0.457080 (0.004971) & 0.698041 (0.010712) &
0.475589 (0.010248) & 0.249147 (0.012639) & 0.524862 (0.014035) & 0.459652\\
\DAWN{} 1.3B &
0.395886 (0.006813) & 0.510157 (0.004989) & 0.723069 (0.010440) &
0.509680 (0.010258) & 0.268771 (0.012955) & 0.507498 (0.014051) & 0.485844\\
\bottomrule
\end{tabular}}
\end{table*}

\subsection{Descriptive larger-scale evaluation at 1.3B}

The 1.272B checkpoint reaches 2.835879 C4 cross-entropy, 17.045382
perplexity, 0.437099 token accuracy, and a 0.485844 six-task mean. It was
trained for 19.999B tokens, about half the 400M training-token count, and no
completed matched dense 1.3B result is admitted to this study. These results
therefore describe a larger trained \DAWN{} model; they are neither a
controlled scaling curve nor a dense-model comparison.

\subsection{Input-conditioned RW operator support}
\label{sec:support}

At all 9,994,240 packed-C4 positions, Q, K, V, and RST each have strict-subset
operator support at positive selection margin (Table~\ref{tab:support}). Mean exact support
is 487.896 Q operators, 386.722 K operators, 812.388 in the shared-pool Q/K
union, 1,649.994 V operators, and 2,017.703 RST operators. This is exact
input-conditioned functional support under Equation~\ref{eq:angular_gate}.

\begin{table}[t]
\centering
\caption{Exact functional support at selection margin \(>0\). Q/K union counts
an operator active for Q or K at the same layer and token.}
\label{tab:support}
\small
\begin{tabular}{lrrr}
\toprule
Computation & Pool & Mean active & Pool fraction\\
\midrule
Q & 4,600 & 487.896 & 10.6064\%\\
K & 4,600 & 386.722 & 8.4070\%\\
Q/K union & 4,600 & 812.388 & 17.6606\%\\
V & 14,500 & 1,649.994 & 11.3793\%\\
RST & 29,300 & 2,017.703 & 6.8864\%\\
\bottomrule
\end{tabular}
\end{table}

The support does not make the present kernel efficient. Static forward
accounting at batch 1 and sequence length 512 is 4.741260116 TFLOPs for
current dense \DAWN{}, 0.818230250 TFLOPs for dense addressing plus
exact-support-only RW application, and 0.386547057 TFLOPs for the matched
dense Transformer. The exact-support opportunity is 17.2577\% of current
\DAWN{} accounting but still 2.1168 times the dense Transformer
(Figure~\ref{fig:flops}). It is a mathematical conditional-RW opportunity,
not measured latency, energy, or hardware speedup.

\begin{figure}[t]
\centering
\includegraphics[width=\linewidth]{figures/flops_stacked.pdf}
\caption{Static backbone accounting at batch 1 and sequence length 512.
Address scoring remains dense in the exact-support estimate. Current execution
also applies all full-dimensional RW rows. No production latency claim is made.}
\label{fig:flops}
\end{figure}

\subsection{IOI circuit localization}

Discovery produces split-half rank stability 0.9139788 over 456,697 ranked
sites. The smallest prefix passing both frozen discovery gates contains 4,096
sites and 0.719249 cumulative absolute importance. It spans all 18 layers and
contains 1,744 Q sites, 2,350 K sites, two V sites, and no RST sites. Split
top-\(k\) overlap is 0.974854 (Jaccard 0.950941). These statistics identify a
Q/K-centered candidate circuit; localization alone is not a causal claim.

\subsection{Validation and held-out causal intervention}
\label{sec:causal}

Table~\ref{tab:ioi} shows confirmatory results. On validation, suppression
moves the mean margin from 3.465456 to \(-4.399331\), a drop of 7.864787
(95\% CI [7.490080, 8.241263]), and accuracy falls from 1.000000 to 0.008130.
The frozen effect exceeds every replicate in all four 100-circuit control
families. The most competitive control is activation-matched, with mean drop
5.400490; the frozen-minus-control difference is 2.464297 (95\% CI
[2.076200, 2.832697], paired permutation \(p=0.00049975\)). Exact restoration
returns accuracy to 0.991870 and recovers 0.999661 of the suppression effect
(95\% CI [0.998655, 1.000584]).

\begin{table*}[t]
\centering
\caption{Frozen IOI circuit results. Accuracy arrows show intact to
suppressed. Recovery is the fraction of the suppression effect recovered by
exact same-example numerator restoration.}
\label{tab:ioi}
\small
\resizebox{0.98\textwidth}{!}{
\begin{tabular}{lrrrrrr}
\toprule
Split & Units & Intact margin & Suppressed margin & Drop (95\% CI) &
Accuracy & Recovery / restored acc.\\
\midrule
Validation & 123 & 3.465456 & -4.399331 &
7.864787 [7.490080, 8.241263] &
1.000000 \(\rightarrow\) 0.008130 & 0.999661 / 0.991870\\
Held-out test & 114 & 3.254720 & -4.155613 &
7.410332 [7.016485, 7.818472] &
1.000000 \(\rightarrow\) 0.008772 & 0.998406 / 1.000000\\
\bottomrule
\end{tabular}
}
\end{table*}

The fully held-out test reuses the unchanged site identity, scoring rule,
intervention, controls, seeds, and thresholds. Suppression moves margin from
3.254720 to \(-4.155613\), a drop of 7.410332 (95\% CI
[7.016485, 7.818472]), while accuracy falls to 0.008772. Again the frozen
effect exceeds every replicate in all four control families. The
activation-matched mean drop is 5.472579; the frozen-minus-control difference
is 1.937754 (95\% CI [1.564818, 2.319706], \(p=0.00049975\)). Restoration
recovers 0.998406 of the effect (95\% CI [0.997324, 0.999385]) and restores
accuracy to 1.000000.

Together, suppression, matched-control separation, and exact restoration
support the claim that this discovery-frozen Q/K-centered operator circuit is
causally necessary for the evaluated IOI behavior on independent validation
and fully held-out inputs. ``Necessary'' is intervention-specific: it does not
mean the sites form a unique or human-readable algorithm.

\subsection{Negative and fail-closed results}
\label{sec:negatives}

The same workflow produces informative negative boundaries
(Table~\ref{tab:negative}). RAVEL passes the behavioral gate but fails
variable-localization stability, so no circuit is frozen and no confirmatory
split is opened. ARC discovery freezes a stable 4,096-site Q/K circuit, but
validation suppression has the wrong sign and fails control separation;
restoration passes, and test remains closed. For 1.3B IOI, behavioral
eligibility and rank stability pass, but the largest allowed 4,096-site prefix
reaches only 0.635144 cumulative importance, below the preregistered 0.70
threshold. No circuit is frozen and no validation or test is opened.

\begin{table*}[t]
\centering
\caption{Fail-closed mechanistic results. Later phases were not opened after a
failed preregistered gate.}
\label{tab:negative}
\small
\begin{tabularx}{\textwidth}{l l Y Y}
\toprule
Study & Passed & Failed boundary & Decision\\
\midrule
RAVEL 400M & Behavioral eligibility & Variable stability:
0.628762 / 0.676243 / 0.660622 & No frozen circuit; no validation, transfer, or test\\
ARC 400M & Stable discovery; 4,096 Q/K sites; restoration &
Suppression effect wrong sign; control separation failed &
Reject causal replication; held-out test closed\\
IOI 1.3B & Behavioral eligibility; rank stability 0.907188 &
4,096-site cumulative importance 0.635144 \(<0.70\) &
No frozen circuit; validation and test closed\\
\bottomrule
\end{tabularx}
\end{table*}

\section{Discussion}
\label{sec:discussion}

\subsection{What the results establish}

The evaluated model shows that a language model can construct attention and
post-attention residual updates by state-conditioned composition of explicit
RW primitives. The matched 400M comparison shows performance
retention under the measured protocols. The support profile shows that
applicability has exact input-conditioned zeros, and the IOI study shows that
operator sites can form a discovery-frozen circuit whose suppression and
restoration causally control behavior on held-out data.

These claims are deliberately narrower than a theory of language reasoning as
discrete programs. The causal result is strong for the specified IOI task and
intervention, not universal across tasks or scales. RAVEL, ARC, and 1.3B IOI
demonstrate that explicit units do not make successful localization automatic.

\subsection{Explicit Transition Construction as a Computational Interface}

The purpose of this decomposition is architectural observability rather than an
assumption of semantic interpretability. A learned primitive may remain
polysemantic, and a transition may depend on a large distributed set of
primitives. \DAWN{} does not claim that such computation is therefore easy for
a human to understand or manipulate. Instead, it asks whether the computation
producing state evolution can first be represented in a form that exposes its
constituent operations, their state-dependent applicability, their executed
contributions, and their composition.

This provides an analysis interface that is absent when the same computation is
represented only as an implicit dense transformation. One may inspect which
primitives participate, suppress or restore their contributions, trace how
applicability changes with the representation state, and test whether localized
compositions are causally involved in behavior. Semantic interpretation of
these objects is a subsequent empirical question rather than a prerequisite of
the architecture.

The central abstraction is a state-dependent map from the current
representation to its applicable RW primitives, their continuously weighted
composition, the resulting composed computation vector, and the state update it
parameterizes:
\[
\begin{array}{c}
\text{current state}\\[-0.15em]
\downarrow\\[-0.15em]
\text{state-conditioned applicability}\\[-0.15em]
\downarrow\\[-0.15em]
\text{RW operator composition}\\[-0.15em]
\downarrow\\[-0.15em]
\text{composed computation vector}\\[-0.15em]
\downarrow\\[-0.15em]
\text{attention or residual update}\\[-0.15em]
\downarrow\\[-0.15em]
\text{next state}
\end{array}
\]
The operator composition makes the intermediate computation explicit. For Q,
K, and V, the composed vectors form the interfaces used by attention; for RST,
the composed vector is applied directly as a residual-state update.

The operation-address space implements one step inside this abstraction. A
query represents the current applicability condition, an address identifies a
reusable primitive for that test, and the RW directions define the primitive's
executed function. Because the address is derived from read/write parameters,
applicability geometry and RW function are coupled rather than assigned
independent labels. This coupling is an architectural constraint, not proof of
semantic interpretability. A primitive may be explicit, observable, and
causally testable while participating in distributed or polysemantic
computation. The IOI circuit's 4,096 sites across all layers is consistent with
this distributed view. Its suppression and exact-restoration results show that
the exposed contributions can serve as causal intervention targets; they do not
show that the circuit is a human-readable algorithm.

\subsection{Toward Low-Dimensional Operation-Space Execution}
\label{sec:future-low-d}

The evaluated model uses low-dimensional addresses but full-\(D\) RW
execution. A separate repository implementation is an \emph{implemented but
not yet fully evaluated low-dimensional extension}. It associates each
operation-space address with \(D\!\leftrightarrow\!R\) interfaces, selects
spaces independently for Q, K, V, and RST, and performs local RW composition
before writing back:
\begin{align}
h_m&=P_m^\top x,\\
z_{a,m}&=\sum_i c_{a,m,i}(h_m)
\langle h_m,r_{a,m,i}\rangle w_{a,m,i},\\
\Delta x_a&=\sum_m\alpha_{a,m}(x)B_m z_{a,m}.
\label{eq:lowd_preview}
\end{align}
Its local operator remains globally rank-1 because
\(B_m w_i r_i^\top P_m^\top=(B_mw_i)(P_mr_i)^\top\). Compilation,
symbolic parameter counts, reduced JIT forward/backward, and gradient
contracts have been checked, and limited early-training observations exist.
There is no completed controlled language-model evaluation or realized
efficiency measurement. The current reference path executes all spaces
densely; candidate-restricted dispatch and an accelerator-native kernel remain
future work. Appendix~\ref{app:lowd} records the boundary in more detail.

\subsection{Limitations}

The current implementation evaluates full address and operator pools densely.
Exact functional support does not imply current sparse execution, and even the
exact-support accounting remains 2.1168 times the matched dense Transformer's
backbone accounting. No production latency, energy, or hardware advantage is
measured.

The central causal evidence is confined to IOI and does not automatically
generalize to every task. RAVEL, ARC, and the selected 1.3B IOI replication
provide valid negative boundaries. The 400M and 1.3B checkpoints have
different parameter and training-token budgets, so they do not define a
controlled scaling curve. Operator-level causal units need not map one-to-one
to human-readable concepts. The low-dimensional extension has not received a
controlled evaluation. Finally, the present work does not establish language
reasoning as universally discrete operator programs.

\section{Conclusion}

\DAWN{} formulates language-model computation as state-conditioned composition
of explicit reusable RW primitives. A representation state determines which RW
primitives are applicable and to what continuous degree; their composition
produces the interfaces used by attention or a direct residual-state update. The
evaluated v4172 realization computes applicability through learned
operation-address mappings derived from RW function and reuses QK-, V-, and
RST-specific primitive pools across depth. It retains measured language-model performance
relative to a matched 400M dense control. A discovery-frozen Q/K-centered
operator circuit is causally necessary under suppression and recoverable under
exact restoration on independent validation and fully held-out IOI inputs. The
model also shows exact input-conditioned functional support, and the frozen IOI
circuit demonstrates that exposed operator contributions can serve as causal
intervention targets. These findings support explicit state-conditioned RW
operator composition as a viable and causally analyzable model of
language-model computation; they do not establish human-readable primitive
semantics, freely editable or transferable computation, benchmark superiority,
or current hardware efficiency.

\appendix

\section{Exact Addressing and Composition Equations}
\label{app:equations}

This appendix restates the production computation as a compact reference. Raw
read and write parameters are normalized in float32 before being cast for
dense contractions:
\[
\bar r_i=\frac{r_i}{\max(\norm{r_i}_2,\epsilon)},\qquad
\bar w_i=\frac{w_i}{\max(\norm{w_i}_2,\epsilon)}.
\]
One pair of model-level probes is shared across QK, V, and RST:
\[
\begin{aligned}
u_i&=\operatorname{unit}(\bar r_i^\top P_r),&
v_i&=\operatorname{unit}(\bar w_i^\top P_w),\\
k_i&=\operatorname{unit}(u_i\od v_i).&&
\end{aligned}
\]
The query and threshold use separate direct projections:
\[
\begin{aligned}
\hat q_a&=\operatorname{unit}(xP_a+b_a),\\
\tau_a&=-1+2\sigma(xu_a+\beta_a),&
\rho_i&=\hat q_a^\top k_i.
\end{aligned}
\]
For \texttt{linear\_angular} composition:
\[
c_i=
\operatorname{clip}\!\left(
\frac{\rho_i-\tau_a}{\max(1-\tau_a,10^{-4})},0,1\right)
\]
and padding rows are masked to zero. This is the single applicability
coefficient used in RW composition. Let \(M_a=\sum_i c_i\). The production
composed output before causal attention or residual addition is
\[
F_a(x)=
\sqrt{\frac{D}{L}}\,
\frac{\sum_i c_i\langle x,\bar r_i\rangle\bar w_i}
{\max(M_a,1)^{p_a}},
\]
with \(p_Q=p_K=0.5\), \(p_V=1.0\), and \(p_{\mathrm{RST}}=1.2\).
Operator rows are sharded over the model axis; local numerators and admission
mass are reduced across shards. Chunking changes the contraction schedule, not
the equation.

\section{Complete Model and Training Settings}
\label{app:settings}

\begin{table*}[h]
\centering
\caption{Detailed settings from the registered training configurations and
restored checkpoint manifests.}
\small
\begin{tabular}{lrrr}
\toprule
Setting & \DAWN{} 400M & Dense 400M & \DAWN{} 1.3B\\
\midrule
Vocabulary & 30,522 & 30,522 & 30,522\\
Residual width & 2,304 & 1,280 & 4,096\\
Layers & 18 & 18 & 24\\
Heads / head width & 36 / 64 & 20 / 64 & 32 / 128\\
Maximum sequence length & 512 & 512 & 512\\
Feed-forward width & -- & 5,120 & --\\
Operation-address width & 256 & -- & 256\\
QK / V / RST pool & 4,600 / 14,500 / 29,300 & -- &
8,534 / 26,904 / 54,352\\
QK / V / RST chunks & 1 / 1 / 4 & -- & 1 / 2 / 8\\
QK / V / RST denominator power & 0.5 / 1.0 / 1.2 & -- &
0.5 / 1.0 / 1.2\\
Initial active targets & 0.05 / 0.03 / 0.01 & -- &
0.05 / 0.03 / 0.01\\
Training batch & 1,024 & 1,024 & 448\\
Mesh (data \(\times\) model) & \(16\times2\) & \(16\times2\) & \(16\times2\)\\
Learning rate & \(3\times10^{-4}\) & \(3\times10^{-4}\) & \(2\times10^{-4}\)\\
Warmup ratio & 0.08 & 0.08 & 0.06\\
Weight decay & 0.1 & 0.1 & 0.1\\
Tau learning-rate multiplier & 0.001 & -- & 0.001\\
Dropout / \texttt{router\_dropout} & 0 / 0 & 0 / -- & 0 / 0\\
Gradient checkpointing & yes & yes & yes\\
Seed & 1 & 1 & 1\\
\bottomrule
\end{tabular}
\end{table*}

The 400M \DAWN{} checkpoint is the primary performance and mechanistic model.
The dense 400M checkpoint is only a matched performance control. The 1.3B
\DAWN{} checkpoint is used for descriptive performance and one fail-closed IOI
localization replication. The still-incomplete dense 1.3B run is excluded.

\section{Benchmark Metrics and Protocol Details}
\label{app:benchmarks}

\smallhead{C4.}
All reported C4 rows use the same packed validation artifact and
\texttt{packed\_c4\_validation\_v1} protocol. The capped view contains 19,531
sequences. Exactly 19,520 sequences form 610 complete global evaluation
batches, producing 9,974,720 valid next-token targets; 11 incomplete-tail
sequences are excluded before scoring.

\smallhead{Zero-shot.}
The protocol is \texttt{mamba\_table3\_zero\_shot\_v1}, using stock
\texttt{lm-eval==0.4.2}, global batch 32, maximum length 512, length buckets
64/128/256/512, and zero few-shot examples. Evaluated split sizes are 5,153
for LAMBADA, 10,042 for HellaSwag, 1,838 for PIQA, 2,376 for ARC-Easy, 1,172
for ARC-Challenge, and 1,267 for WinoGrande. Task, dataset, tokenizer, and
software contracts match exactly between the two 400M models.

\section{IOI Localization and Intervention Details}
\label{app:ioi}

\subsection{Frozen discovery}

Discovery admits 119 paired-correct independent units and ranks 456,697 sites.
Split-half rank stability is 0.913978765. The frozen selection rule is the
smallest prefix satisfying cumulative absolute importance at least 0.70 and
split top-\(k\) overlap at least 0.95. The selected \(k=4096\) prefix has:
\begin{itemize}
    \item cumulative absolute importance 0.719248713;
    \item 3,993 shared split sites, overlap 0.974853516, and Jaccard
    0.950940700;
    \item Q/K/V/RST counts \(=1744/2350/2/0\);
    \item coverage of every layer from 0 through 17.
\end{itemize}
Validation and test do not recompute or alter this selection.

\begin{table*}[!t]
\centering
\caption{Control effects and frozen-minus-control separation. Each control
family contains 100 replicates. All paired permutation \(p\)-values are
0.0004997501 with 2,000 permutations.}
\label{tab:controls}
\small
\begin{tabular}{lrrrr}
\toprule
Control & Validation drop & Frozen minus control (95\% CI) &
Test drop & Frozen minus control (95\% CI)\\
\midrule
Equal-count & 0.070886 & 7.793902 [7.437851, 8.161572] &
0.080276 & 7.330057 [6.922408, 7.716879]\\
Same-layer & 0.039389 & 7.825399 [7.440629, 8.180288] &
0.041073 & 7.369260 [6.977447, 7.754539]\\
Activation-matched & 5.400490 & 2.464297 [2.076200, 2.832697] &
5.472579 & 1.937754 [1.564818, 2.319706]\\
Route-frequency & 0.219162 & 7.645626 [7.289512, 8.031625] &
0.242665 & 7.167668 [6.779836, 7.545537]\\
\bottomrule
\end{tabular}
\end{table*}

\subsection{Suppression, controls, and restoration}

For a correct candidate \(c\) and source candidate \(s\), the primary score is
\[
\mathcal{M}=\sum_{t\in c}\log p(t)-\sum_{t\in s}\log p(t).
\]
Suppression zeros the frozen sites' execution numerator at specified token
positions. It leaves admission mass and the production denominator unchanged.
Controls use the same 4,096-site count and exclude all circuit sites.
Activation matching conditions on layer, computation type, and the discovery mean
pre-scale operator-output-norm quantile. Restoration executes a distinct
intervention: after suppression, it adds back the exact selected same-example
numerator under the unchanged denominator.

On validation, restoration returns the margin to 3.462791, leaving residual
drop 0.002665. On held-out test, it returns the margin to 3.242905, leaving
residual drop 0.011815. The respective recovery fractions are 0.999661155
(95\% CI [0.998654544, 1.000584125]) and 0.998405645
(95\% CI [0.997323737, 0.999385266]).

\section{Negative Mechanistic Results}
\label{app:negative}

\smallhead{RAVEL.}
Discovery uses 142 independent units. Rank stability is 0.628762 for
Continent, 0.676243 for Country, and 0.660622 for Language, below the
preregistered requirement. No prefix passes the complete discovery rule, no
variable circuit is frozen, and validation, transfer, and test are not opened.

\smallhead{ARC.}
Discovery stability is 0.89386 and a 4,096-site Q/K-only circuit is frozen.
On 51 validation units, intact margin is 1.735534 and suppressed margin is
2.819139, giving a frozen margin drop of \(-1.083605\) (95\% CI
[\(-1.472180\), \(-0.690348\)]). The effect has the wrong sign and fails
control separation. Restoration recovery is 1.004083 (95\% CI
[0.994587, 1.013652]), but restoration alone does not pass the conjunction.
Held-out test remains closed.

\smallhead{1.3B IOI.}
The scale replication admits 125 independent units and passes rank stability
at 0.907187749. The largest allowed 4,096-site prefix reaches only 0.635144
cumulative absolute importance, below the 0.70 gate. No circuit is frozen;
validation and test remain closed.

\section{Compute Accounting}
\label{app:compute}

The support profile evaluates 9,994,240 packed-C4 positions and observes
address-stage margins before RW application. It stores aggregate statistics,
not per-token operator identities or vectors. Static accounting uses batch 1,
sequence length 512, and \(1\) MAC \(=2\) FLOPs. It counts major
matmul/einsum operations and omits softmax, normalization, gate elementwise
arithmetic, embedding lookup, residual additions, and communication
collectives.

\begin{table*}[!t]
\centering
\caption{Backbone static accounting. ``Exact support'' retains full dense
address scoring and replaces only full-pool RW application with the measured
positive-margin support.}
\small
\begin{tabular}{lrrr}
\toprule
Component & Current dense \DAWN{} & Dense address + exact RW & Dense Transformer\\
\midrule
Address scoring + query/tau projections & 0.293742 & 0.293742 & --\\
RW application & 4.306187 & 0.383157 & --\\
Causal attention & 0.043487 & 0.043487 & 0.024159\\
Dense output / QKV / FFN maps & 0.097845 & 0.097845 & 0.362388\\
\midrule
Total TFLOPs & 4.741260 & 0.818230 & 0.386547\\
\bottomrule
\end{tabular}
\end{table*}

The exact-support estimate is mathematically output-equivalent if the executor
can skip all and only zero-weight RW terms. It does not include indexed
address retrieval and does not demonstrate kernel parity. Instrumented
profiling time is excluded from latency claims because histogram collection is
part of the measured path.

\section{Low-Dimensional Operation-Space Extension}
\label{app:lowd}

The low-dimensional extension uses one shared atlas of \(M=8\) operation-space
addresses, shared \(D\!\to\!R\) read interfaces \(P_m\), and shared
\(R\!\to\!D\) write interfaces \(B_m\). Q, K, V, and RST have separate query
projections, tau maps, and operator banks, and independently select two atlas
entries. For local operator \(i\) in space \(m\),
\[
\begin{aligned}
\op^{\mathrm{global}}_{a,m,i}(x)
&=B_m w_{a,m,i}r_{a,m,i}^{\top}P_m^\top x\\
&=(B_mw_{a,m,i})(P_mr_{a,m,i})^\top x,
\end{aligned}
\]
so the induced global transformation remains rank-1.

The 40M reference has 40,341,088 parameters, \(D=432\), \(R=128\), 12
layers, six heads, and top-2 space selection. Repository-local compilation,
abstract-tree and symbolic parameter counting, shared-key and independent
Q/K/V/RST space selection, legacy-checkpoint rejection, reduced JIT
forward/backward, selector metrics, and gradient contracts have passed.
Some early-training observations are recorded, but the exact deployment and
complete controlled language-model evaluation are absent. No latency,
throughput, energy, or hardware-efficiency result exists. The reference path
executes every space densely; future work requires candidate-restricted
dispatch and an accelerator-native execution kernel.

\section{Reproducibility Identities}
\label{app:repro}

\begin{table*}[!t]
\centering
\caption{Primary immutable identities. Branch names are omitted because they
move; exact checkpoint steps, hashes, and artifact locations are the binding
identities.}
\scriptsize
\begin{tabularx}{\textwidth}{>{\raggedright\arraybackslash}p{0.13\textwidth}
>{\raggedright\arraybackslash}p{0.43\textwidth} Y}
\toprule
Evidence & Exact identity & Repository source\\
\midrule
\DAWN{} 400M & model \texttt{spatial-r1-v4.1.7.2}; step 76,293;
checkpoint identity \texttt{a7ce8afc...f91c17}; config hash
\texttt{08733ae4...5039b} &
\path{configs/train_config_v4172_400M_c4_40B_v4_64_ver1_den_qk0p5_v1p0_rst1p2.yaml}\\
Dense 400M & model \texttt{baseline-JAX}; step 76,293; config hash
\texttt{c3ce5d72...82fc25} &
\path{configs/train_config_baseline_tpu_400M_c4_40B_v4_64_tp2.yaml}\\
\DAWN{} 1.3B & model \texttt{spatial-r1-v4.1.7.2}; step 87,193;
checkpoint identity \texttt{2e38b41f...bf5f9}; config hash
\texttt{1e877263...4a5222} &
\path{configs/train_config_v4172_1p3B_c4_20B_v4_64_ver1_den_qk0p5_v1p0_rst1p2.yaml}\\
IOI circuit & parsed frozen-spec hash \texttt{6ec26618...418c};
site hash \texttt{573cfff8...0200}; selected-row hash
\texttt{1439f6f3...fe0c1} &
\path{configs/paper_ioi_frozen_circuit_v4172_400m.yaml}\\
Compute support & support hash \texttt{b1eac5f5...bdc2};
accounting hash \texttt{d7ad6115...984e} &
\path{assets/paper/v4172_400m_compute_support_step76293/}\\
\bottomrule
\end{tabularx}
\end{table*}

The machine-readable checkpoint registry is
\texttt{configs/paper\_checkpoint\_registry.yaml}. Exact evaluation commits,
software versions, data hashes, result roots, and complete artifact hashes are
recorded in the repository experiment ledger. Historical experiment branches
are evidence identities only; this source is maintained on the canonical main
branch.

\section*{Acknowledgments}

The author thanks Anthropic Claude, Claude Code, OpenAI ChatGPT, and OpenAI
Codex for assistance throughout this research. This work was developed with
extensive use of AI systems for literature exploration, conceptual exploration,
cross-domain framing, abstraction and terminology refinement, implementation
support, debugging, and manuscript preparation. In particular, these systems
were used as interactive assistants for refining the learned state-transition
formulation and the exposition of \DAWN{}. The research direction,
architectural decisions, experimental design, experiments, analyses,
interpretations, and final claims remain the responsibility of the author.

This research was supported by the Google TPU Research Cloud (TRC) program,
which provided computational resources for training and analysis.

\begingroup
\small
\setlength{\bibsep}{0pt}
\bibliography{references}
\bibliographystyle{icml2026}
\endgroup

\end{document}
